problem:  
Let $\Sigma^k$ be an exotic smooth sphere of dimension $k \ge 7$.  
Denote by $\mathcal{R}_{\ge0}(\Sigma^k)$ the space of all smooth Riemannian metrics on $\Sigma^k$ with **nonnegative sectional curvature**, equipped with the $C^\infty$ topology, and let  
\[
\mathcal{M}_{\ge0}(\Sigma^k) = \mathcal{R}_{\ge0}(\Sigma^k) / \mathrm{Diff}_0(\Sigma^k)
\]
be the corresponding **moduli space of nonnegatively curved metrics modulo diffeomorphisms isotopic to the identity**.  
Similarly define $\mathcal{M}_{\ge0}(S^k)$ for the standard sphere.

**Problem:**  
Does there exist a dimension $k \ge 7$ and an exotic sphere $\Sigma^k$ such that $\mathcal{M}_{\ge0}(\Sigma^k)$ and $\mathcal{M}_{\ge0}(S^k)$ are *not homotopy equivalent*?  
Equivalently, can the exotic smooth structure on $\Sigma^k$ be detected by any topological invariant of the moduli space of nonnegatively curved metrics (for example, the number of connected components, $\pi_1$, or higher homotopy groups)?  

In particular, is it possible that $\mathcal{M}_{\ge0}(\Sigma^k)$ has more connected components than $\mathcal{M}_{\ge0}(S^k)$, reflecting geometrically distinct classes of nonnegatively curved metrics arising from the exotic differential structure?  
Or, conversely, is the smooth exoticity of $\Sigma^k$ completely invisible within the topology of its curvature moduli space?  

Why is it a "good" problem:  
This version is sharply formulated and mathematically sound: it fixes a precise moduli space $\mathcal{M}_{\ge0}$ (quotienting by $\mathrm{Diff}_0$) and asks whether exotic smooth structures can alter its fundamental topological features. It directly probes the **detectability of smooth structure through curvature‑restricted moduli spaces**, a question that lies squarely at the frontier of current Riemannian and differential topology.  

Resolving it would require deep insight into how nonnegative curvature interacts with the topology of diffeomorphism groups and smoothing theory. A positive answer would reveal that exotic smoothness leaves measurable traces in geometric moduli, while a negative answer would imply a surprising rigidity of nonnegative curvature under smooth variation. Either outcome would open new territory between **infinite‑dimensional moduli theory**, **curvature geometry**, and **surgery‑theoretic classification of manifolds**, unifying ideas from these distinct domains.  

The problem is also commendably simple to state—comparing two moduli spaces associated with homeomorphic manifolds—but is profound in depth: any progress would require breakthroughs in our understanding of the topology of spaces of nonnegatively curved metrics, making it a genuinely “good” and forward‑looking research problem.